\subsection{Independent valuations}

DEFINITION 1. Let A and A' be two valuation rings of the same field K. A and A' are called independent if K is the ring generated by A and A'. Two valuations on K are called independent if their rings are independent and dependent otherwise.

An improper valuation on K is independent of every valuation on K. For two valuations of height 1 on K to be independent, it is necessary and sufficient that they be not equivalent ( \(\S\) 4, no. 5, Proposition 6 (c)).

THEOREM 1 (Approximation Theorem for Valuations). Let \(v_{i}\) ( \(1 \leqslant i \leqslant n\) ) be valuations on a field \(\mathbf{K}\) which are independent in pairs and \(\Gamma_{i}\) the order group of \(v_{i}\) . Let \(a_{i} \in \mathbf{K}\) and \(\alpha_{i} \in \Gamma_{i}\) ( \(1 \leqslant i \leqslant n\) ). Then there exists \(x \in \mathbf{K}\) such that \(v_{i}(x - a_{i}) \geqslant \alpha_{i}\) for all \(i\) .

If \(v_{i}\) is improper, then \(\alpha_{i} = 0\) and the relation \(v_{i}(x - a_{i}) \geqslant \alpha_{i}\) is true for all \(x \in \mathbf{K}\) . We may therefore assume that the \(v_{i}\) are not improper.

Let \(\mathbf{A}_i\) be the ring of \(v_i\) , \(\mathbf{B} = \bigcap_{i=1} \mathbf{A}_i\) and \(\mathfrak{p}_i = \mathfrak{m}(\mathbf{A}_i) \cap \mathbf{B}\) . By Proposition 1 of no. 1, the \(a_i\) may be written \(a_i = b_i / s\) ( \(b_i \in \mathbf{B}\) , \(s \in \mathbf{B} - \{0\}\) ); if we write \(x = y / s\) and \(\alpha_i' = \alpha_i + v_i(s)\) , then \(v_i(y - b_i) \geqslant \alpha_i'\) . This shows that we may assume that \(a_i \in \mathbf{B}\) for all \(i\) ; we may also assume that \(\alpha_i > 0\) for all \(i\) . Let \(\mathfrak{v}_i\) be the set of \(z \in K\) such that \(v_{i}(z) \geqslant \alpha_{i}\) ; we write \(q_{i} = v_{i} \cap B\) . For \(x \in B\) , \(v_{i}(x - a_{i}) \geqslant \alpha_{i}\) is equivalent to \(x \equiv a_{i}(q_{i})\) . We therefore need to show that the canonical homomorphism

\(\mathbf{B} \to \prod_{i=1}^{n} \left( \mathbf{B} / q_i \right)\) is surjective, that is that \(q_i + q_j = \mathbf{B}\) for \(i \neq j\) (Chapter II, § 1, no. 2, Proposition 5). As the maximal ideals of \(\mathbf{B}\) are the \(p_i\) (Proposition 2), it will suffice for this to show that \(q_i \not\subset p_j\) for \(i \neq j\) .

Suppose that there exists \(i, j\) such that \(\mathfrak{q}_i \subset \mathfrak{p}_j\) and \(i \neq j\) . We shall see shortly that the radical of \(\mathfrak{q}_i\) is a prime ideal \(\mathfrak{p}\) of B. Then \(\mathfrak{p} \subset \mathfrak{p}_j\) and also \(\mathfrak{p} \subset \mathfrak{p}_i\) since \(\alpha_i > 0\) and hence \(\mathfrak{q}_i \subset \mathfrak{p}_i\) . Therefore \(\mathrm{A}_j = \mathrm{B}_{\mathfrak{p}_j} \subset \mathrm{B}_{\mathfrak{p}}\) (no. 1, Proposition 1) and similarly \(\mathrm{A}_i \subset \mathrm{B}_{\mathfrak{p}}\) . Now, as \(\mathfrak{v}_i \neq (0)\) and \(\mathfrak{v}_i = \mathrm{B}_{\mathfrak{p}_i} \mathfrak{q}_i\) (Chapter II, §2, no. 4, Proposition 10), \(\mathfrak{q}_i \neq (0)\) , whence \(\mathfrak{p} \neq (0)\) and \(\mathrm{B}_{\mathfrak{p}} \neq \mathrm{K}\) . This contradicts the hypothesis that \(\mathrm{A}_i\) and \(\mathrm{A}_j\) are independent.

It remains to show that p is prime. Now this follows from the following lemma:

LEMMA 2. Let A be a valuation ring and b an ideal of A distinct from A. Then the radical r of b is a prime ideal.

Suppose that \(xy \in \mathfrak{r}\) . Then there exists \(n \geqslant 1\) such that \((xy)^n \in \mathfrak{b}\) . Let \(v\) be a valuation associated with A. If, for example, \(v(x) \geqslant v(y)\) , then

\[
v (x ^ {2 n}) \geqslant v (x ^ {n} y ^ {n}),
\]

whence \(x^{2n} \in \mathfrak{b}\) and \(x \in \mathfrak{r}\) .

COROLLARY 1. For every family of elements \(\gamma_{i} \in \Gamma_{i}\) ( \(1 \leqslant i \leqslant n\) ), there exists \(x \in \mathbf{K}\) such that \(v_{i}(x) = \gamma_{i}\) ( \(1 \leqslant i \leqslant n\) ).

We may assume that \(\mathbf{A}_i \neq \mathbf{K}\) for all \(i\) . Then, there exists for all \(i\) an \(a_i \in \mathbf{K}\) such that \(v_i(a_i) = \gamma_i\) and an \(\alpha_i \in \Gamma_i\) such that \(\gamma_i < \alpha_i\) . We apply Theorem 1 to these elements \(a_i\) : there exists \(x \in \mathbf{K}\) such that \(v_i(x - a_i) > v_i(a_i)\) ; whence, as \(x = a_i + (x - a_i)\) , \(v_i(x) = v_i(a_i) = \gamma_i\) ( \(\S 3\) , no. 1, Proposition 1).

COROLLARY 2. Let \(\mathcal{T}_i\) be the topology defined on \(\mathbf{K}\) by \(v_i\) ; let \(\mathbf{K}^n\) be given the topology the product of the \(\mathcal{T}_i\) . If the \(v_i\) are not improper, the diagonal of \(\mathbf{K}^n\) is dense in \(\mathbf{K}^n\) .

PROPOSITION 3. Let \(v\) and \(v'\) be two non-improper valuations on the same field \(K\) . For \(v\) and \(v'\) to define the same topology on \(K\) , it is necessary and sufficient that they be dependent.

Suppose that the topologies \(\mathcal{T}_v\) and \(\mathcal{T}_{v'}\) , defined by \(v\) and \(v'\) , are identical. Since \(\mathcal{T}_v\) is Hausdorff, the diagonal of \(\mathbf{K}^2\) is closed and hence \(v\) and \(v'\) are dependent (Corollary 2 to Theorem 1).

Conversely, suppose that \(v\) and \(v'\) are dependent. Then their rings A and A' are contained in the same ring A'' distinct from K and A'' is the ring of a valuation \(v''\) (§ 4, no. 1, Proposition 1). It suffices to show that the topology \(\mathcal{T}_{v''}\) is identical with \(\mathcal{T}_v\) . Let \(\Gamma\) and \(\Gamma''\) be the order groups of \(v\) and \(v''\) . There exists an increasing homomorphism \(\lambda\) of \(\Gamma\) onto \(\Gamma''\) such that \(v'' = \lambda \circ v\) (§ 4, no. 3). If \(\alpha'' \in \Gamma''\) , let \(\alpha \in \overline{\lambda}^{1}(\alpha'')\) ; the condition \(v(x) \geqslant \alpha\) implies \(v''(x) \geqslant \alpha''\) . Let \(\beta \in \Gamma\) and \(\beta'' = \lambda(\beta)\) ; the condition \(v(x) \leqslant \beta\) implies \(v''(x) \leqslant \beta''\) and hence the condition \(v''(x) > \beta''\) implies \(v(x) > \beta\) . As \(v\) and \(v''\) are not improper, the inequalities in question define fundamental systems of neighbourhoods of 0 for \(\mathcal{T}_v\) and \(\mathcal{T}_{v''}\) . Hence \(\mathcal{T}_v = \mathcal{T}_{v''}\) , which completes the proof.

Remarks

\begin{enumerate}
\def\labelenumi{(\arabic{enumi})}
\item
  Proposition 3 shows that the relation ``v and v' are dependent'' is an equivalence relation.
\item
  Taking account of the relations between valuations of height 1 and ultrametric absolute values (§ 6, no. 2), Proposition 3 also follows, in the case of valuations of height 1, from the characterization of equivalent absolute values (General Topology, Chapter IX, § 3, no. 2, Proposition 5).
\end{enumerate}

PROPOSITION 4. Let \(v_{1}, \ldots, v_{n} (n \geqslant 2)\) be pairwise dependent valuations on the same field K. Then the rings \(A_{1}, \ldots, A_{n}\) of \(v_{1}, \ldots, v_{n}\) generate a subring of K distinct from K.

For \(n = 2\) Proposition 4 follows from Definition 1. Suppose it holds for \(n - 1\) valuations. Then there exists a subring A of K distinct from K and containing \(\mathbf{A}_1, \ldots, \mathbf{A}_{n-1}\) ; there also exists a subring \(\mathbf{B} \neq \mathbf{K}\) containing \(\mathbf{A}_{n-1}\) and \(\mathbf{A}_n\) . As A and B contain \(\mathbf{A}_{n-1}\) , they are comparable with respect to inclusion (§ 4, no. 1, Corollary to Proposition 1). The greater of these two therefore contains all the \(\mathbf{A}_i\) .

